\chapter*{ВИСНОВКИ}
\addcontentsline{toc}{chapter}{ВИСНОВКИ}
\thispagestyle{fancy}

У дипломній роботі розроблено VPN-застосунок з двофакторною автентифікацією
на основі протоколу WireGuard. За результатами роботи можна зробити такі висновки:

\begin{enumerate}
  \item Проаналізовано існуючі технології VPN та встановлено переваги протоколу
        WireGuard з точки зору продуктивності, безпеки та простоти реалізації.

  \item Досліджено архітектуру протоколу WireGuard та його криптографічні основи
        (Curve25519, ChaCha20-Poly1305, BLAKE2s, Noise~IK).

  \item Досліджено методи двофакторної автентифікації та обґрунтовано вибір
        стандарту TOTP як оптимального з огляду на простоту інтеграції та
        доступність клієнтських застосунків.

  \item Проведено аналіз існуючих рішень та виявлено відсутність комплексного
        академічно обґрунтованого підходу до інтеграції 2FA з WireGuard.

  \item Розроблено архітектуру та реалізовано програмний комплекс, що поєднує
        WireGuard з двофакторною автентифікацією.

  \item Проведено тестування розробленої системи та підтверджено її функціональність
        і відповідність поставленим вимогам.

  \item Виконано порівняльний аналіз продуктивності розробленої системи відносно
        чистого ядерного WireGuard та OpenVPN на однакових VPS-вузлах (Ubuntu~26.04
        LTS, UpCloud PL-WAW1 ↔ DE-FRA1, RTT~23{,}94~мс). Система продемонструвала
        паритет затримки (0\,\% за tcp\_p, +0{,}66~мс під навантаженням) та
        пропускної здатності TCP (різниця 1{,}0\,\% за tcp\_t, 0\,\% за tcp\_p)
        порівняно з чистим WireGuard; час введення у тунель становив 1571~мс
        проти 92~мс для чистого \texttt{wg-quick up}.
\end{enumerate}

Підтверджено наукову новизну запропонованого підходу~--- інтеграція 2FA з
WireGuard без модифікації протоколу та без використання попередньо
узгоджених ключів.

\newpage
